\noindent
This paper, the eleventh in a series on the philosophy of intimacy and the theory of justice, develops the empirical praxis of the good relation whose conditions the prior paper established in the abstract. Where the prior paper held that the practice of a good relation is the cultivation of a field rather than the construction of a structure, that the matter of this field is the everyday sensory and habitual fabric of a shared life, and that a good field satisfies two conditions, the eudaimonic and the just, of which the just constitutes rather than merely constrains the eudaimonic, the present paper asks how a shared life is in fact cognised, practised, assessed, and sustained across its course. It treats that course as a cycle, and the cycle as its primary object: a recurrence of cognition, practice, evaluation, reflection, and renewed cognition, turning across the time of a relation from its beginning, through its marriage, into its later years. The cycle is shown to be a real cycle in the dynamical sense, not a linear sequence of stages but a self-sustaining movement whose persistence is the condition of the relation's existence, so that the cessation of the cycle is not a relation at rest but a relation extinguished.

Five frameworks for the cognition and practice of flourishing are distinguished, the positivist, the evidentialist, the normative, the virtue-theoretic, and the phenomenological, and they are treated not as competing theories of one measurable object but as the symmetry-broken phases, in the sense of the prior paper, of an object no single framework possesses. The frameworks are the resources the cycle draws upon at the moments suited to each, and the paper proceeds by traversing the moments of the cycle in turn, drawing the appropriate frameworks into each. A cognitive and computational basis for the whole is established first: appraisal, before it is an explicit measurement, is a continuous and largely unconscious inference upon the partially observable state of a relation, conducted in a value that is not symbolised, and formalised as the belief updating of a decision process whose objective is itself revised by its running, which is the decision-theoretic form of the prior paper's thesis that there is no fixed meta-value-grammar. On the free-energy reading of this inference, the continual appraisal by which a relation is maintained is the activity by which it holds itself in existence, and the cessation of appraisal is the beginning of its extinction.

The paper carries three commitments through its length. The first is an operational concession: a practice requires what is assessable, and the paper therefore restricts itself, where it must assess, to value as structured into an assessable proxy, in the explicit awareness that the proxy is not the value the prior papers held to be transversal across the registers and irreducible to symbolisation. This is the first occasion in the series on which an incomplete value is adopted deliberately, and it is the honest discharge of the prior result rather than its abandonment. The second commitment concerns the assessment of justice: a relation's distribution of returned value, taken relative to a baseline set by the coupling of the parties, is shown to admit a statistical signature, the skewness of the distribution, by which the exploitation of one party within an apparently satisfying relation becomes visible where aggregate measures of satisfaction conceal it. The justice of a relation's value cycle is rendered, under the concession, as a testable hypothesis about the symmetry of a return distribution against a coupling-based baseline, modelled as a branching process on the relational network. The third commitment concerns the cycle as a whole: the cycle accumulates a phase, in the sense of the prior paper's holonomy, and its regime is fixed by two independent questions, whether it persists and, if it persists, what sign of phase it accumulates. Three regimes follow, the flourishing spiral of a persisting cycle of positive phase, the catastrophic cycle of a persisting cycle of negative phase, and the extinction of a cycle that ceases to turn, with the passage between them understood as a bifurcation. Justice is shown to be the internal condition of the flourishing spiral and not a constraint upon it, since a relation whose return distribution is severely skewed cannot accumulate a positive phase for the party consigned to its low-return tail.

In keeping with the thesis that no single framework possesses the whole, the paper terminates not in a synthesis but in plural conclusions, each stated to the boundary of its framework, and it holds throughout that the flourishing of a shared life is cultivated and not constructed, assessed in proxy and not possessed, and sustained as an open spiral and not secured as a state.
